\chapter{Series numéricas}
\label{cap:series-numericas}

Una serie se estudia mediante la sucesión de sus sumas parciales. El capítulo ordena criterios de convergencia, distingue convergencia absoluta de condicional y practica la justificación de cada prueba \cite{apostol1974,abbott2015,rudin1976}.

\section{Series y sumas parciales}
Definición de serie, convergencia y valor como límite de sumas parciales.

\section{Condición necesaria de convergencia}
El término general debe tender a cero; ejemplos que muestran que no basta.

\section{Series de términos positivos}
Monotonía de las sumas parciales y criterios básicos de convergencia.

\section{Criterios de comparación}
Comparación directa y por límite con series de referencia.

\section{Criterios del cociente y de la raíz}
Pruebas para series positivas y tratamiento de los casos inconclusos.

\section{Convergencia absoluta}
Convergencia absoluta, reordenamientos y comparación con convergencia simple.

\section{Series alternantes}
Criterio de Leibniz y estimación del resto.

\section{Convergencia condicional}
Ejemplos de convergencia condicional y cautela con los reordenamientos.
